\documentclass[11pt]{article}
\usepackage[T1]{fontenc}
\usepackage[margin=1in]{geometry}
\usepackage{amsmath,amssymb,amsthm,mathtools}
\usepackage[colorlinks=true,linkcolor=blue,urlcolor=blue,citecolor=blue]{hyperref}
\usepackage{microtype}
\newtheorem{theorem}{Theorem}[section]
\newtheorem{lemma}[theorem]{Lemma}
\newtheorem{corollary}[theorem]{Corollary}
\newtheorem{proposition}[theorem]{Proposition}
\theoremstyle{remark}
\newtheorem{remark}[theorem]{Remark}
\newcommand{\one}{\mathbf 1}
\newcommand{\E}{\mathbb E}
\newcommand{\Tr}{\operatorname{Tr}}
\newcommand{\dd}{\,\mathrm d}
\newcommand{\SU}{\mathrm{SU}(2)}
\newcommand{\norm}[1]{\left\lVert#1\right\rVert}
\setlength{\emergencystretch}{3em}
\hypersetup{pdftitle={Marginal boundary energy costs in the Wilson model},pdfauthor={Research note},pdfsubject={Finite lattice estimates, not a continuum Yang--Mills construction}}
\title{Marginal boundary energy costs in the Wilson model}
\author{Research note}
\date{October 9, 2026}
\begin{document}
\maketitle
\begin{abstract}
The preceding rough-region note disproved all-small filling with arbitrary
frozen boundary links. We now pay for large boundary-forced energy in the
\emph{unconditioned} Wilson measure, instead of assuming it can always be
removed. On four-dimensional periodic cubes of side $L=2^m$, $m\ge2$,
site-reflection positivity, the standard event chessboard estimate, and a
64-class link decomposition give an exponential moment bound for the energy
of every fixed edited-link set. The same bound controls its conditional
minimum in the actual exterior marginal. It also gives a counted tail for
source-defined rough-link components, without a boundary compatibility
hypothesis. Constants and overlap losses are explicit. These estimates do
not produce an all-small compatible filling, an unrestricted lower bound
on removable energy, or control of general renormalized or reweighted measures. A sublevel kernel additionally gives partial cooling on a low-minimum
class, with a separately paid exceptional label in the raw measure.
These are finite Wilson-model inputs, not a continuum construction, mass-gap proof, or novelty
claim. The imported chessboard theorem and the locally proved geometric
application are distinguished below.
\end{abstract}
\tableofcontents

\section{Model and main estimates}
Work on $\Lambda_L=(\mathbb Z/L\mathbb Z)^4$, $L=2^m$, $m\ge2$.
There are $V=L^4$ sites, $4V$ positively oriented links and $6V$ plaquettes.
Each link has a variable $U_e\in\SU$, the reverse link has $U_e^{-1}$,
and the reference measure is product \emph{probability} Haar. Set
\[
 s(g)=1-\tfrac12\Tr g=\tfrac12\norm{\one-g}_{\rm op}^2,
 \quad E(U)=\sum_p s(U_p),\quad
 \mu_H(\dd U)=Z_H^{-1}e^{-HE(U)}\prod_e\dd U_e,\quad H\ge1.
\]
For a fixed link set $J$, let $n=|J|$ and let $P(J)$ contain \emph{all}
plaquettes touching $J$. Write
\[
 E_J(U)=\sum_{p\in P(J)}s(U_p),\qquad
 e_{*,J}(z)=\min_{x\in\SU^J}E_J(x,z),\quad z=U_{J^c}.
\]
The minimum exists and is continuous in $z$. It obeys
$0\le e_{*,J}(U_{J^c})\le E_J(U)$. Define
\begin{equation}
 a_H=\tfrac32+\log24+\tfrac32\log(8H),\qquad
 b_H=\tfrac12a_H+\frac{\log2}{64}.
 \label{eq:constants}
\end{equation}
These constants are deliberately loose and differ from the second pass's
compatible-kernel constants.

\begin{theorem}[Unconditioned local energy and boundary cost]
\label{thm:moment}
For every fixed $J$ and $0<\theta<1$,
\begin{equation}
 \E_{\mu_H}\exp\!\left(\frac{\theta H}{64}E_J\right)
 \le \exp\!\left[n\left(\theta a_H-
                  \frac{\log(1-\theta)}{64}\right)\right].
 \label{eq:moment}
\end{equation}
The same bound holds with $E_J$ replaced by
$e_{*,J}(U_{J^c})$. In particular, for $v\ge0$,
\begin{align}
 \mu_H(E_J\ge v)&\le\min\{1,e^{-Hv/128+b_Hn}\},\label{eq:energytail}\\
 \mu_H(e_{*,J}(U_{J^c})\ge v)
 &\le\min\{1,e^{-Hv/128+b_Hn}\}.\label{eq:mintail}
\end{align}
The second probability is under the Wilson \emph{exterior marginal}.
It is not a statement conditional on an arbitrary fixed exterior.
\end{theorem}

For a roughness parameter $p>0$, call a plaquette rough when
$\norm{\one-U_p}_{\rm op}>p/\sqrt H$, and write $N_p(J)$ for the number
of rough plaquettes in $P(J)$. The letter $p$ here is a parameter, not a
plaquette index. All such events are empty if $p/\sqrt H\ge2$.

\begin{corollary}[Rough counts without boundary compatibility]
\label{cor:rough}
For fixed $J$ and a positive integer $k$,
\begin{equation}
 \mu_H(N_p(J)\ge k)
 \le\min\{1,\exp[-kp^2/256+b_Hn]\}.
 \label{eq:rough}
\end{equation}
\end{corollary}
\begin{proof}
The event implies $E_J>kp^2/(2H)$. Apply \eqref{eq:energytail}.
\end{proof}
These conclusions do not say that the rough plaquettes can be removed
while preserving $J^c$. Their original energy and their boundary-forced
minimum are separate quantities.

\section{A volume-proportional lower bound on the partition function}
Let $C_\rho$ be the geodesic cap in $\SU$ of angle at most $\rho$, with
$0<\rho\le1$. The previous note's elementary Haar calculation gives
\[
 \operatorname{Haar}(C_\rho)
 =\frac{\rho-\sin\rho\cos\rho}{\pi}\ge\frac{\rho^3}{24}.
\]
If every link lies in this cap, each plaquette has operator curvature at
most $4\rho$, by telescoping its four unitary factors, and hence
$E\le48V\rho^2$. Take $\rho=(8H)^{-1/2}$. Integrating just over this
product cap gives
\begin{equation}
 Z_H\ge(\rho^3/24)^{4V}e^{-48HV\rho^2}
       =e^{-4Va_H}.
 \label{eq:partition}
\end{equation}
No independence of plaquette matrices, gauge fixing, or weak-coupling
expansion is used. Probability Haar and the nonnegative-action convention
are essential to the subsequent numerator bound by one.

\section{Site reflection and the chessboard input}
\label{sec:reflection}
Reflect a coordinate $x_i$ to $2a-x_i$ modulo $L$, where $a$ is an integer.
The reflection has two fixed planes, at $a$ and $a+L/2$. An oriented link
is transported geometrically; if its orientation reverses, its matrix is
inverted. This is a Haar-preserving involution. Tangential links on the
fixed planes remain pointwise fixed. Plaquette energy is preserved,
because reversal and change of base point invert or conjugate holonomy.

Let $z$ be the links wholly contained in these two planes, and $x_+,x_-$
the remaining links in the two closed half-lattices. No elementary
plaquette involves strict interiors on both sides: its extent in the
reflected coordinate is zero or one, while the reflecting planes have
integer coordinates. Consequently
\[
 E=E_P(z)+E_+(x_+,z)+E_-(x_-,z),
 \qquad E_-(x_-,z)=E_+(\vartheta x_-,z).
\]
For a bounded real function $f$ supported in the positive closed half,
\begin{equation}
 \E_{\mu_H}[f\,\vartheta f]
 =\frac1{Z_H}\int e^{-HE_P(z)}
   \left[\int e^{-HE_+(x_+,z)}f(x_+,z)\dd x_+\right]^2\dd z\ge0.
 \label{eq:rp}
\end{equation}
Reflection symmetry also gives the symmetric bilinear form. This proves
site-reflection positivity directly. It does not claim positivity for an
arbitrary boundary-conditioned law or for an arbitrary extra interaction.

\paragraph{Imported standard theorem.}
Tile by closed cubes of side two,
$B_t=a+2t+\{0,1,2\}^4$, indexed by
$\mathcal T=(\mathbb Z/(L/2)\mathbb Z)^4$, with
$M=|\mathcal T|=V/16$; the origin $a$ can be translated.
Neighboring blocks share boundary links. For an event $A$ on the base
block, let $\vartheta_t A$ denote its reflected translate to block $t$:
local coordinate $v_j$ becomes $v_j$ for even $t_j$ and $2-v_j$ for odd
$t_j$, followed by translation by $a+2t$.
The event chessboard inequality is
\begin{equation}
 \mu_H\!\left(\bigcap_{j=1}^r\vartheta_{t_j}A_j\right)
 \le\prod_{j=1}^r
 \mu_H\!\left(\bigcap_{t\in\mathcal T}\vartheta_t A_j\right)^{1/M},
 \qquad t_1,\ldots,t_r\text{ distinct}.
 \label{eq:chessboard}
\end{equation}
We use the event formulation in Biskup--Koteck\'y
\cite[\S2.2, Theorem 2.4, (2.10)]{bk}, retaining its power-of-two
restriction on $L/2$. Its usual repeated Cauchy--Schwarz proof applies to
the closed-half \emph{edge algebras} just described: the needed facts are
support under reflection, multiplication, commuting coordinate
reflections, and \eqref{eq:rp}. These hold for links at their geometric
midpoints, including inversion of reversed orientations. Shared-face
variables belong to both closed halves and are integrated as $z$ in
\eqref{eq:rp}; they are not independently copied. Thus neither independent
blocks nor an identification of outgoing links with spins at their tails
is assumed. This is the standard chessboard argument on an edge-decorated
lattice, not a new chessboard theorem.

Only the event version \eqref{eq:chessboard} is needed below. We derive
exponential moments from it explicitly, rather than silently importing a
stronger functional version. The geometric extension to these edge
algebras and the ensuing analysis are written arguments, not Lean
certificates or external expert verification.

\section{Reflected stars and the 64 link classes}
For a link $e$, write
\[
 E_e=E_{\{e\}}=\sum_{q\ni e}s(U_q).
\]
A prototype link in a side-two block has direction $\mu$, local starting
coordinate $r\in\{0,1\}$ in direction $\mu$, and starting coordinate one
in each of the other three directions. All six incident plaquettes lie
in the block. Each of their relative interiors lies inside that block:
two coordinates span a unit interval, and the other two equal one.
In particular no such plaquette is wholly on a block face.

It follows that the reflected copies of a prototype star have
\emph{disjoint plaquette sets}. They may share links on block faces;
no probabilistic independence follows or is needed. Their energies obey
\begin{equation}
 \sum_{t\in\mathcal T}E_{e(t)}\le E.
 \label{eq:disjoint}
\end{equation}
For $u\ge0$, dissemination of $\{E_e\ge u\}$ forces $E\ge Mu$.
Its constrained partition numerator is at most $e^{-HMu}$ since total
reference mass is one. Equations \eqref{eq:partition} and
\eqref{eq:chessboard} therefore give
\begin{equation}
 \mu_H\!\left(\bigcap_t\{E_{e(t)}\ge u\}\right)^{1/M}
 \le q_H(u):=\min\{1,\exp[-Hu+64a_H]\}.
 \label{eq:starweight}
\end{equation}
The factor $64$ here is $4V/M=4\cdot16$, not the number of plaquettes
in a star.

Partition all links by their direction and the four parities of their
starting site: $4\cdot2^4=64$ classes. Fix direction $\mu$ and parity
vector $c$. Choose the block grid with
$a_\mu=c_\mu$ and $a_\nu\equiv c_\nu-1\pmod2$ for $\nu\ne\mu$.
Every link of this class is the lower prototype ($r=0$) in a distinct
block of that grid, and all of its incident plaquettes are in that block.
When pulling an event back to the base block, odd $t_\mu$ replaces this
prototype by $r=1$; both prototypes satisfy \eqref{eq:disjoint}.
Thus \eqref{eq:chessboard} allows different pulled-back prototypes.

\begin{lemma}[Joint star tails within a class]\label{lem:classtails}
For a fixed set $S$ in one of these 64 classes and thresholds $u_e\ge0$,
\begin{equation}
 \mu_H(E_e\ge u_e\text{ for all }e\in S)
 \le\prod_{e\in S}q_H(u_e).
 \label{eq:classtails}
\end{equation}
The same bound holds with each $E_e$ replaced by
$e_{*,\{e\}}(U_{e^c})$, by event inclusion.
\end{lemma}
\begin{proof}
The events occupy distinct blocks of the common grid. Each reflected
prototype has the bound \eqref{eq:starweight}. Apply
\eqref{eq:chessboard}. For the final assertion use
$e_{*,\{e\}}\le E_e$ pointwise.
\end{proof}
In particular the sharper one-link boundary bound is
\begin{equation}
 \mu_H(e_{*,\{e\}}\ge u)\le q_H(u).
 \label{eq:singleminimum}
\end{equation}
For the second pass's staple mismatch
$\mathcal D=\sum_{i<j}\norm{S_i-S_j}_{\rm op}^2$,
$e_{*,\{e\}}=\mathcal D/(6+|\sum s_j|)\ge\mathcal D/12$, hence
\begin{equation}
 \mu_H(\mathcal D\ge d)\le\min\{1,e^{-Hd/12+64a_H}\}.
 \label{eq:mismatch}
\end{equation}
This pays a marginal statistical cost for large one-link frozen-loop
mismatch. It is not an estimate conditional on that mismatch having
already been frozen.

\section{From joint tails to arbitrary edited sets}
Put $K=64a_H$ and $T_e=HE_e$. Within a class,
\eqref{eq:classtails} bounds joint upper tails by products of
$q(v)=\min\{1,e^{K-v}\}$. For $0<\theta<1$ use
\[
 e^{\theta T_e}=1+\int_0^\infty
                  \theta e^{\theta v}\one_{\{T_e>v\}}\dd v.
\]
Expand a finite product and apply Tonelli and
Lemma~\ref{lem:classtails} to every subset. The result is
\begin{align}
 \E\exp\!\left(\theta\sum_{e\in S}T_e\right)
 &\le\left[1+\int_0^\infty\theta e^{\theta v}q(v)\dd v\right]^{|S|}
   \nonumber\\
 &=\left[\frac{e^{\theta K}}{1-\theta}\right]^{|S|}.
 \label{eq:classmoment}
\end{align}
The equality follows by splitting the integral at $K\ge0$.
This is an upper-tail comparison, not a claim of independent energies.

\begin{proof}[Proof of Theorem~\ref{thm:moment}]
For $J$, write $S_c$ for its intersections with the 64 classes. Every
plaquette in $P(J)$ is counted at least once in $\sum_{e\in J}E_e$, so
$E_J\le\sum_{e\in J}E_e$. Apply H\"older with 64 equal exponents and
then \eqref{eq:classmoment}:
\begin{align*}
 \E e^{\theta HE_J/64}
 &\le \E\prod_c\exp\!\left(\frac\theta{64}
                                      \sum_{e\in S_c}T_e\right)\\
 &\le\prod_c\left[\E\exp\!\left(\theta
                                      \sum_{e\in S_c}T_e\right)\right]^{1/64}\\
 &\le\exp\!\left[\frac n{64}
                   (\theta K-\log(1-\theta))\right].
\end{align*}
Empty classes contribute one; no independence across classes is used.
Substitute $K=64a_H$. The minimum bound follows pointwise from
$e_{*,J}\le E_J$. Markov's inequality at $\theta=1/2$ proves the tails.
\end{proof}

\begin{corollary}[Explicit overlap loss for retained minima]
\label{cor:overlap}
Let $J_1,\ldots,J_r$ be fixed nonempty link sets, let $I=\bigcup_iJ_i$,
and let
\[
 m_* =\max_{q\in P(I)}|\{i:q\in P(J_i)\}|.
\]
Then
\begin{equation}
 \E\exp\!\left[\frac{H}{128m_*}
                  \sum_i e_{*,J_i}(U_{J_i^c})\right]
 \le e^{b_H|I|}.
 \label{eq:overlap}
\end{equation}
For disjoint link sets, $m_*\le4$, not necessarily one.
\end{corollary}
\begin{proof}
Pointwise, $\sum_i e_{*,J_i}\le\sum_i E_{J_i}\le m_*E_I$.
Apply \eqref{eq:moment}. A plaquette has four distinct links, so it
can touch at most four disjoint $J_i$.
\end{proof}
This controls a sum of retained energy labels. It does not make energy
gains from simultaneous replacements additive.

\section{A near-minimum kernel for every boundary}
The marginal bound can be paired with a replacement that retains, rather
than deletes, the unavoidable minimum. This is a probability-kernel
construction, not an efficient sampling algorithm.

\begin{proposition}[Uniform sublevel replacement]\label{prop:sublevel}
For every fixed nonempty $J$ there is a Borel gauge-covariant kernel,
independent of the original edited links, which freezes $J^c$, has
conditional Haar density at most $(24H^{3/2})^n$, and obeys for every output
\begin{equation}
 E_J(y,z)\le e_{*,J}(z)+12n/H,\qquad
 N_p(J;y,z)\le \frac{2He_{*,J}(z)+24n}{p^2}.
 \label{eq:sublevel}
\end{equation}
Its edited-link law depends only on the frozen links of affected plaquettes.
It need not remove every rough plaquette or preserve reflection positivity.
\end{proposition}
\begin{proof}
Define
\[
 B_J(z)=\{y:E_J(y,z)\le e_{*,J}(z)+12n/H\},\quad
 v_J(z)=\operatorname{Haar}_J(B_J(z)),\quad
 K_z(\dd y)=\frac{\one_{B_J(z)}(y)}{v_J(z)}\dd y.
\]
To lower-bound $v_J$, choose any minimizer $V$ for this one fixed $z$.
For $y_e=V_e e^{X_e}$ with $\norm{X_e}_{\rm op}\le\rho=H^{-1/2}$, the
second pass's minimizer-cap estimate gives
$E_J(y,z)\le e_{*,J}(z)+12\sum_e\norm{X_e}^2$.
For clarity, its proof differentiates $E_J(Ve^{tX},z)$ twice: each
plaquette contributes at most $(\sum_{e\in\partial q\cap J}\norm{X_e})^2$,
and four links per plaquette and six plaquettes per link give
$|E_J''(t)|\le24\sum_e\norm{X_e}^2$. The first derivative is zero at
an unrestricted minimizer. Taylor's integral remainder supplies the
factor 12. Thus $B_J(z)$ contains a product cap of Haar mass at least
$(\rho^3/24)^n$, proving the density bound.

Continuity of $E_J$ and $e_{*,J}$ makes the joint sublevel set Borel;
parameter integration makes $v_J$ Borel and strictly positive. This proves
measurability without choosing a minimizer measurably. A site gauge maps
$B_J(z)$ bijectively to $B_J(z^g)$ and preserves edited-link Haar, proving
covariance and equality of normalizing volumes. Copying the exterior
leaves all unaffected plaquettes unchanged. Finally each rough output
plaquette contributes more than $p^2/(2H)$ energy, yielding
\eqref{eq:sublevel}.
\end{proof}

\begin{corollary}[Partial cooling with a paid exceptional label]
\label{cor:partialcooling}
For fixed $J$ and positive integer $k$, put
\[
 A=\{N_p(J)\ge k,\ e_{*,J}\le kp^2/(4H)\},\qquad
 \ell_H=\log24+\tfrac32\log H.
\]
If $kp^2\ge96n$, the kernel of Proposition~\ref{prop:sublevel} sends
sources in $A$ to at most $3k/4$ rough affected plaquettes, with
\begin{equation}
 H(E(x)-E(y))\ge kp^2/8.
 \label{eq:partialgain}
\end{equation}
If also $kp^2\ge16n\ell_H$, then for every nonnegative measurable
$\mathcal Q$ with $0<Z_{\mathcal Q}=\int e^{-HE}\mathcal Q<\infty$,
\begin{equation}
 I_{\mathcal Q}(A):=\frac1{Z_{\mathcal Q}}
       \int_A e^{-HE(x)}(K\mathcal Q)(x)\dd x
 \le e^{-kp^2/16}.
 \label{eq:partialintegral}
\end{equation}
The complementary high-minimum label has the \emph{raw Wilson} bound
\begin{equation}
 \mu_H(e_{*,J}>kp^2/(4H))
 \le\min\{1,e^{-kp^2/512+b_Hn}\}.
 \label{eq:exception}
\end{equation}
If $kp^2\ge1024b_Hn$, the last expression is at most
$e^{-kp^2/1024}$.
\end{corollary}
\begin{proof}
On $A$, \eqref{eq:sublevel} gives
$N_p(J;y,z)\le k/2+24n/p^2\le3k/4$. The source energy is greater than
$kp^2/(2H)$ and output energy at most $kp^2/(4H)+12n/H$, giving
$H(E(x)-E(y))\ge kp^2/4-12n\ge kp^2/8$.

For completeness, the bounded-density replacement argument is as follows.
Writing $x=(\xi,z)$ and $y=(\eta,z)$, use the pointwise energy drop on
$A$ to replace $e^{-HE(\xi,z)}$ by
$e^{-kp^2/8}e^{-HE(\eta,z)}$. Bound the edited-output density by
$M_H=e^{n\ell_H}$; integrating the original $\xi$ over its probability
Haar space costs at most one. Tonelli leaves exactly
$Z_{\mathcal Q}$ as the remaining integral. Hence
$I_{\mathcal Q}(A)\le e^{-kp^2/8+n\ell_H}\le e^{-kp^2/16}$.
Apply \eqref{eq:mintail} for \eqref{eq:exception} and absorb $b_Hn$
under the last numerical budget.
\end{proof}
The literal integral in \eqref{eq:partialintegral} evaluates $\mathcal Q$
at the \emph{output}; it is not the $\mathcal Q$-weighted probability of
$A$. For $\mathcal Q=1$ it equals the raw probability. The exceptional
label in \eqref{eq:exception} is not paid for an arbitrary $\mathcal Q$.
Consequently these two inequalities cannot be combined into an
unrestricted weighted/RG replacement theorem. Some affected plaquettes
may become newly rough: the conclusion concerns their \emph{total count},
not containment in the old rough inventory. Adaptive supports still
require the following counting step, and iteration needs gains and density
bounds along the actual sequence of outputs.

\section{Counting source-dependent rough regions}
Connect two lattice links when they share a plaquette. This graph has
maximum degree 18. A connected set of $n$ links containing a given root
has a deterministic spanning-tree traversal of length $2(n-1)$, and the
traversal determines its visited set. Hence there are at most
$18^{2(n-1)}$ such sets.

Let $\mathcal R_p(U)$ be \emph{all} links incident to at least one rough
plaquette. Its connected components are determined by the source
configuration. For a root link $e_0$, let $C_p(e_0)$ be its component,
with size zero if $e_0\notin\mathcal R_p$. Put
\begin{equation}
 \kappa_H(p)=p^2/1024-b_H.
 \label{eq:kappa}
\end{equation}

\begin{theorem}[Actual rough-component size tail]\label{thm:components}
For every integer $n\ge1$,
\begin{equation}
 \mu_H(|C_p(e_0)|\ge n)
 \le\min\{1,18^{2(n-1)}e^{-\kappa_H(p)n}\}.
 \label{eq:componenttail}
\end{equation}
Thus there is an exponentially decreasing size bound when
$\kappa_H(p)>2\log18$. For $p=(\log H)^P$ and any $P>1/2$, this condition
holds for all sufficiently large $H$.
\end{theorem}
\begin{proof}
A component of size at least $n$ contains a connected $n$-link subset $J$
through $e_0$. Every link in $J$ has a rough witness plaquette in $P(J)$.
One plaquette can witness at most four links, so
$N_p(J)\ge\lceil n/4\rceil$. For each \emph{fixed} such $J$,
Corollary~\ref{cor:rough} bounds this event by
$\exp[-np^2/1024+b_Hn]$. Sum over the deterministic list of possible
$J$. This treats adaptive selection by a union bound, not by conditioning
on the event that a particular $J$ was selected. Finally
$b_H=(3/4)\log H+O(1)$ proves the last assertion.
\end{proof}

There is a similar bound for connected retained-minimum labels. If
$\lambda=Hu/128-b_H>2\log18$ and $N\ge1$, then
\begin{equation}
 \mu_H\!\left(\exists J\ni e_0\text{ connected},\ |J|\ge N,
          \ e_{*,J}\ge u|J|\right)
 \le\min\left\{1,
 \frac{18^{2(N-1)}e^{-\lambda N}}{1-18^2e^{-\lambda}}\right\}.
 \label{eq:adaptivemin}
\end{equation}
This follows by summing \eqref{eq:mintail} over geometries and sizes.
No regular conditional Gibbs formula is applied after choosing $J$.

\section{What this supplies and what it does not}
\paragraph{The supplied input.}
Large boundary-forced energy now has a volume-uniform \emph{marginal}
cost in the original finite Wilson model. The one-link holonomy mismatch
has the explicit tail \eqref{eq:mismatch}. The components of the actual
rough-link set have a paid source-selection count and the witness density
$N_p(J)\ge |J|/4$ needed for \eqref{eq:componenttail}. Arbitrary overlapping
fixed supports can be handled with the declared loss \eqref{eq:overlap}.
None of these statements assumes the second pass's forest compatibility.
Proposition~\ref{prop:sublevel} also supplies an every-boundary near-minimum
kernel. Corollary~\ref{cor:partialcooling} gives a genuine energy gain and
rough-count reduction on its specified low-minimum class, with an explicit
raw-measure cost for the complementary label.

\paragraph{Not enough removable energy.}
Outside that low-minimum class, the quantity available to a boundary-preserving edit is at most
$E_J-e_{*,J}$. Bounds on the separate marginal tails of $E_J$ and $e_{*,J}$
do not bound their ratio or difference on a selected rough event.
The frustrated star still has $E_e=e_{*,\{e\}}=6$ for every central-link
value when its exterior is fixed. Its large-minimum exterior is expensive
in \eqref{eq:singleminimum}, but its conditional removal gain is still zero.
One cannot replace marginal rarity by conditional repairability.
Likewise one isolated $-\one$ link in identity background has total
energy 12, regardless of padding the editable region: it cannot pay a
volume-proportional replacement entropy in an arbitrarily enlarged set.

\paragraph{A minimum is not a full holonomy label.}
Take constant commuting links $U_0=\exp(i\phi\sigma_3/L)$, the other
directions identity, with $0<\phi<\pi$. Every plaquette is flat, so
$e_{*,J}=0$ for every $J$. Choose $J$ to be the direction-one links
starting at $(x_0,0,0,0)$ for all $x_0$. Its frozen touched graph contains
the full direction-zero loop at $x_1=x_2=x_3=0$, of holonomy
$\exp(i\phi\sigma_3)\ne\one$. Thus even zero boundary-forced plaquette
energy does not enforce trivial frozen cycles for arbitrary, possibly
wrapping, regions. Small-energy estimates cannot silently replace the
specified forest compatibility condition.

\paragraph{The raw measure is essential.}
For $\mathcal Q=e^{HE}$, reweighting $\mu_H$ gives product Haar. For one
plaquette its holonomy is Haar, since its four links are distinct and
independent. Therefore
$\mu_{H,\mathcal Q}(s(U_q)>1)=1/2$ for every $H$, whereas the corresponding
raw Wilson tail tends to zero. General positive weights, signed RG
expansions, or extra long-range interactions do not inherit this result.
Their reflection positivity or a replacement/comparison estimate would
have to be proved separately. The bound also does not cover arbitrary
fixed boundary conditions, non-dyadic volumes, or all Gibbs states.

\paragraph{Next analytic task.}
Either strengthen the partial-cooling dichotomy to an all-small,
boundary-compatible repair with controlled enlargement and a lower bound on
$E_J-e_{*,J}$ per edited link, or build a retained-defect expansion whose
weights, normalizers, and interaction changes remain controlled after
renormalization. A new kernel must still track all affected plaquettes,
its Haar density, output-evaluated weights, and overlap along the actual
sequence of edits. Neither task is accomplished by the marginal bounds
here. The fixed-physical-box defect test, cutoff comparison, nontrivial
continuum fields and reconstruction remain downstream obligations.

\section{Verification and literature scope}
The accompanying script checks the 64-class embedding, reflected-star
plaquette disjointness, the absence of plaquettes crossing both strict
halves, SU(2) reflection invariance, the constants and tail integrations,
and the obstruction, padding, overlap and flat-holonomy negative controls.
It also checks event chessboard inequalities by exact enumeration in a
small two-dimensional $\mathbb Z_2$ gauge model. That finite-group model is
a diagnostic of the bookkeeping, not a numerical evaluation of the
four-dimensional $\SU$ partition function or a proof of its probabilities.
No new analytic result is Lean-formalized; the earlier scalar certificates
retain their original scope. This pass is self-audited, not independently
agent-reviewed or externally certified.

The only imported analytic inequality is the standard chessboard estimate
in \cite{bk}. Older gauge-theory literature already discusses
volume-uniform suppression of large plaquette fluctuations; see, for
context only, \cite[\S1.1, (4)]{bvf} and the Mack--Petkova paper cited
there. We do not import that reported bound or its asymptotic constants,
and make no novelty claim for the present deliberately weak application.
Exact source and access details are recorded in \texttt{sources.md}.

\begin{thebibliography}{9}
\bibitem{bk}
M.~Biskup and R.~Koteck\'y,
\emph{Forbidden gap argument for phase transitions proved by means of
chessboard estimates}, Commun. Math. Phys. \textbf{264} (2006), 631--656,
arXiv:math-ph/0505011. Author manuscript, \S2.2, Definition 2.2 and
Theorem 2.4, (2.10), pp.~4--5.
\bibitem{bvf}
O.~Borisenko, S.~Voloshin and M.~Faber,
\emph{Plaquette representation for 3D lattice gauge models:
I. Formulation and perturbation theory}, arXiv:hep-lat/0508003,
\S1.1, (4), and reference [20]. Context only, not a proof input.
\end{thebibliography}
\end{document}
